% =============================================================================
% Kapitel 2 -- Anforderungen und Pattern-Auswahl
% =============================================================================
\section{Anforderungen und Pattern-Auswahl}

Das Spiel erfüllt alle funktionalen Anforderungen: drei Turmtypen (Scout, Sniper,
Rapid), drei Gegnertypen (Standard, Fast, Tank) mit unterschiedlichen
Geschwindigkeiten und Trefferpunkten, ein vollständiger Rundenzyklus mit vier
Phasen sowie zehn verschiedene Kartentypen für Bau, Buffs und Sonderaktionen.
Grafisch dargestellt wird das Spiel über libGDX mit Spielfeld, HUD und
Kartenhand. Nicht-funktional wurde auf Modularität und Erweiterbarkeit geachtet:
Abhängigkeiten verlaufen über Interfaces, neue Karten- oder Turmtypen lassen sich
ohne Änderung bestehenden Codes ergänzen. Das Spielsystem
wurde in Java~21 mit libGDX umgesetzt und per Gradle gebaut;
Lombok reduziert Boilerplate durch \texttt{@Getter}-Annotationen.

Die Auswahl der neun GoF-Patterns folgt direkt aus den Anforderungen -- jedes
Muster löst ein identifizierbares Entwurfsproblem. Tabelle~\ref{tab:patterns}
gibt einen Überblick.

\begin{table}[H]
\centering
\begin{tabular}{llp{7.5cm}}
\toprule
\textbf{Pattern} & \textbf{Kategorie} & \textbf{Gelöstes Problem} \\
\midrule
Singleton      & Erzeugung & Exakt eine konsistente Spielzustandsinstanz \\
Factory Method & Erzeugung & Deck kennt keine konkreten Kartentypen \\
Prototype      & Erzeugung & Konfigurierte Objekte klonen statt neu aufbauen \\
Builder        & Erzeugung & Turmkonfiguration vor Platzierung trennen \\
Decorator      & Struktur  & Türme zur Laufzeit stapelbar erweitern \\
Command        & Verhalten & Kartenaktionen als eigenständige Objekte kapseln \\
Strategy       & Verhalten & Zielwahl und Angriffstyp unabhängig kombinieren \\
State          & Verhalten & Phasenverhalten ohne if-else-Ketten verwalten \\
Observer       & Verhalten & Ereignisse entkoppelt von reaktiven Systemen \\
\bottomrule
\end{tabular}
\caption{Übersicht der neun implementierten GoF-Design-Patterns}
\label{tab:patterns}
\end{table}


Besonders hervorzuheben ist die Aufteilung in vier Erzeugungsmuster, die gemeinsam
die gesamte Objekterzeugung des Spiels abdecken, und vier Verhaltensmuster, die
das laufende Spielgeschehen strukturieren. Das Decorator-Pattern als einziges
Strukturmuster nimmt dabei eine Sonderrolle ein: Es ist das architektonisch
prägendste Muster des Projekts und kooperiert eng mit Command und Strategy.
